\headline{{\textbf{\Large\sffamily Seasonal Care Tips:}}\\
{\textbf{\large\sffamily Mid\--December to Mid\--January}}}
\label{2024-12-care}
\noindent
Winter might seem like a quiet time in the bonsai calendar, but mid\--December to mid\--January offers plenty of opportunities to care for your trees and prepare for the year ahead.
With a little attention, you can keep your bonsai healthy and thriving through the colder months.

\topic{1. Protect Your Trees from Winter Weather}
\noindent
London winters are unpredictable --- frosty mornings can quickly give way to heavy rain.
While hardy species like pines and maples can handle cold weather, they still need some care.

\begin{itemize}
    \item \textbf{Frost Protection:} Move trees to a sheltered spot or wrap smaller pots in fleece or bubble wrap if temperatures drop below -2°C.

    \item \textbf{Manage Excess Rain:} Ensure pots are draining well to avoid waterlogged roots.
    If needed, tilt pots slightly or move them under shelter during prolonged rain.
\end{itemize}

\topic{2. Prune and Tidy}
\noindent
Winter is ideal for light pruning and general tidying.

\begin{itemize}
    \item \textbf{Structure Assessment:} With deciduous trees bare, you can clearly see their branch structure.
    Remove any deadwood, crossed branches, or twigs that disrupt the design.
    Major pruning, however, should wait until late winter.

    \item \textbf{Care for Evergreens:} For pines and junipers, thin out old or browning needles to improve air circulation, but avoid over\--thinning.
\end{itemize}

\columnbreak
\topic{3. Monitor Watering}
\noindent
Watering in winter can be tricky.
Trees need much less water when dormant, but their roots can still dry out, especially in windy weather.

\begin{itemize}
    \item \textbf{Check Soil Regularly:} Water lightly if the topsoil feels dry, preferably in the morning to avoid overnight freezing.

    \item \textbf{Avoid Waterlogging:} Ensure pots have proper drainage to prevent root rot.
\end{itemize}

\topic{4. Watch for Pests and Diseases}
\noindent
Winter doesn't completely eliminate pests or fungal problems.

\begin{itemize}
    \item \textbf{Inspect Regularly:} Look for scale insects, aphids, or spider mites in bark crevices.
    Remove them with a cotton swab and rubbing alcohol.

    \item \textbf{Prevent Fungal Issues:} Ensure good airflow around your trees and consider a fungicide if necessary.
\end{itemize}

\topic{5. Prepare for Spring}
\noindent
Though it's not yet time for repotting, this is the moment to plan ahead.
Stock up on soil, pots, and tools, and check for roots growing out of drainage holes --- a sign that repotting may be needed soon.

\topic{6. Enjoy Winter Bonsai}
\noindent
Winter's bare branches reveal the beauty of your trees' structure, while evergreens stand out against the muted landscape.
Take time to appreciate this unique season, perhaps photographing your trees or creating a simple indoor display.

With a bit of care and preparation, your bonsai will be ready to thrive when spring arrives.
Enjoy the season, and happy growing!
\closearticle
%\clearpage
